% SLICEDEPTH2 research note
\documentclass[11pt]{article}
\usepackage[margin=.9in]{geometry}
\usepackage{amsmath,amssymb,amsthm}
\newtheorem{theorem}{Theorem}
\newtheorem{lemma}{Lemma}
\begin{document}

\begin{center}
{\Large SLICEDEPTH2: Exact depth-two lower bounds for sparse-slice 3-UNSAT}\\[4pt]
\textbf{Date: 2026-10-04}\\
\textbf{Status: CANDIDATE; symbolic proof complete, independent review pending}\\
\textbf{P versus NP status: OPEN}
\end{center}

\section{Setup}

Let
\[
M=8\binom q3,\qquad t=6q,
\]
and let \(S_{q,t}\) be the monotone slice extension of clause-presence
\(\mathrm{UNSAT}_{3,q}\):
\[
S_{q,t}(A)=0\ (|A|<t),\qquad
S_{q,t}(A)=\mathrm{UNSAT}_{3,q}(A)\ (|A|=t),\qquad
S_{q,t}(A)=1\ (|A|>t).
\]

SLICEPAD1 proves, subject to independent review, that the exact-\(t\) SAT
language is NP-complete and the exact-\(t\) UNSAT language is coNP-complete.

\section{Minterms and maxterms}

Every UNSAT \(A\in\binom{\mathcal C_q}{t}\) is a minterm of \(S_{q,t}\):
every strict subset has weight below \(t\) and is therefore a zero-input.

Every SAT \(B\in\binom{\mathcal C_q}{t}\) is a maxterm of \(S_{q,t}\):
every strict superset has weight above \(t\) and is therefore a one-input.

For a monotone DNF computing a monotone function, each minterm forces a
distinct term equal to that minterm.  For a monotone CNF, each maxterm forces
a distinct clause equal to the complement of that maxterm.

\begin{lemma}[DNF minterm forcing]
If a monotone DNF computes \(S_{q,t}\), then it has at least as many terms as
there are UNSAT members of \(\binom{\mathcal C_q}{t}\).
\end{lemma}

\begin{proof}
Let \(A\) be a minterm.  Some conjunction term \(T\) is satisfied by \(A\),
so its support is contained in \(A\).  The support of \(T\) itself is a
one-input, because setting exactly those variables to one satisfies the term.
Minimality of \(A\) therefore forces \(T=A\).  Distinct minterms force
distinct terms.
\end{proof}

\begin{lemma}[CNF maxterm forcing]
If a monotone CNF computes \(S_{q,t}\), then it has at least as many clauses
as there are SAT members of \(\binom{\mathcal C_q}{t}\).
\end{lemma}

\begin{proof}
Let \(B\) be a maxterm.  Some positive clause \(K\) is falsified by \(B\), so
\(K\subseteq\mathcal C_q\setminus B\).  For every coordinate
\(j\notin B\), the strict superset \(B\cup\{j\}\) is a one-input and must
satisfy \(K\).  Hence \(j\in K\).  Thus
\(K=\mathcal C_q\setminus B\).  Distinct maxterms force distinct clauses.
\end{proof}

\section{Counting UNSAT minterms}

For a fixed assignment \(\alpha\), exactly \(7M/8\) clause coordinates are
satisfied by \(\alpha\).  Therefore
\[
\#\{\text{SAT }t\text{-sets}\}
\le
2^q\binom{7M/8}{t}.
\]
Consequently
\[
\#\{\text{UNSAT }t\text{-sets}\}
\ge
\binom Mt
-
2^q\binom{7M/8}{t}.
\]
Moreover
\[
\frac{\binom{7M/8}{t}}{\binom Mt}
=
\prod_{i=0}^{t-1}
\frac{7M/8-i}{M-i}
\le
\left(\frac78\right)^t,
\]
so with \(t=6q\),
\[
\#\{\text{UNSAT }t\text{-sets}\}
\ge
\left[
1-\left(2(7/8)^6\right)^q
\right]\binom M{6q}.
\]
Since
\[
\rho:=2(7/8)^6=0.89759063720703125<1,
\]
the bracket tends to one exponentially fast.

\section{Counting SAT maxterms}

Fix one assignment \(\alpha\).  Every \(t\)-subset of the \(7M/8\) clauses
satisfied by \(\alpha\) is satisfiable.  Therefore
\[
\#\{\text{SAT }t\text{-sets}\}
\ge
\binom{7M/8}{6q}.
\]

\section{Depth-two theorem}

\begin{theorem}[SLICE-DEPTH2]
For all sufficiently large \(q\), every monotone DNF computing
\(S_{q,6q}\) has at least
\[
\left(1-\rho^q\right)\binom M{6q}
\]
terms, and every monotone CNF computing \(S_{q,6q}\) has at least
\[
\binom{7M/8}{6q}
\]
clauses, where \(M=8\binom q3\) and
\(\rho=2(7/8)^6<1\).

In particular, both sizes are
\[
\exp(\Omega(q\log q))
=
M^{\Omega(M^{1/3})}.
\]
\end{theorem}

\begin{proof}
The exact counting bounds follow from the preceding lemmas.

For the asymptotic estimate, use
\[
\binom Nk\ge\left(\frac Nk\right)^k.
\]
Since \(M=\Theta(q^3)\) and \(k=6q\),
\[
\left(\frac{cM}{6q}\right)^{6q}
=
\left(\Theta(q^2)\right)^{6q}
=
\exp(\Theta(q\log q))
\]
for any fixed \(c>0\).  Also \(q=\Theta(M^{1/3})\), giving the stated form.
For DNF, \(1-\rho^q\) is bounded below by a positive constant for all
sufficiently large \(q\).
\end{proof}

\section{Interpretation}

This is a genuine unconditional lower bound inside a restricted circuit
model, conditional only on the exact function definition, not on \(P\ne NP\)
or ETH.  It is nevertheless far from the project target.

Depth-two circuits have no multi-level sharing.  The open problem is to show
that allowing arbitrary monotone DAG sharing cannot compress the exact-layer
transversal family to polynomial size.  Wegener's slice examples show that
minterm/maxterm abundance alone is insufficient for this.

The SLICEKW1 note identifies the correct size model as rectangle-DAG /
DAG-like monotone Karchmer--Wigderson protocols.  The present theorem supplies
a baseline adversarial requirement for any proposed DAG decomposition:
it must compress exponentially many exact-layer minterms and maxterms through
shared feasible rectangles.

\section{Next target}

Search for a structural obstruction to such sharing.  Candidate parameters
should be tested on the exact block system
\[
B_\alpha=\{C:\alpha\text{ falsifies }C\},
\qquad
|B_\alpha\cap B_\beta|
=
\binom{q-d_H(\alpha,\beta)}3.
\]
The target must distinguish this hard slice from Wegener-style easy slices
and from Horn, 2-SAT, affine/Tseitin, and bounded-treewidth controls.

\section{Status}

\[
\boxed{
\text{depth-two monotone lower bound: CANDIDATE with complete proof}
}
\]
\[
\boxed{
\text{general monotone circuit lower bound: OPEN}
}
\]
Thus \(P\ne NP\) remains OPEN.

\end{document}
